\documentclass[aspectratio=169,11pt]{beamer}

% --- Theme & Color Scheme ---
\usetheme{Madrid}
\usecolortheme{whale}

% --- Custom Neutral/Professional Colors ---
\definecolor{workstationdark}{RGB}{24, 24, 24}
\definecolor{workstationblue}{RGB}{0, 122, 204}
\definecolor{softgray}{RGB}{245, 245, 245}
\definecolor{accentcyan}{RGB}{56, 189, 248}

\setbeamercolor{palette primary}{bg=workstationdark,fg=white}
\setbeamercolor{palette secondary}{bg=workstationblue,fg=white}
\setbeamercolor{palette tertiary}{bg=workstationdark,fg=white}
\setbeamercolor{structure}{fg=workstationblue}
\setbeamercolor{title}{bg=workstationdark,fg=white}
\setbeamercolor{frametitle}{bg=workstationdark,fg=white}
\setbeamercolor{block title}{bg=workstationblue,fg=white}
\setbeamercolor{block body}{bg=softgray,fg=black}

% Remove default navigation symbols
\setbeamertemplate{navigation symbols}{}

% Custom slide numbers in footer
\setbeamertemplate{footline}[frame number]

% --- Packages ---
\usepackage{amsmath,amssymb}
\usepackage{booktabs}
\usepackage{graphicx}

% --- Title Information ---
\title{\textbf{Direct Volume Rendering \& Surface Display}}
\subtitle{Architectural Synthesis of Marc Levoy (1988) \& Robert Drebin et al. (1988)}
\author{\textbf{Tarun P Krishnamurthy} \and \textbf{Harsh Gunda} \and \textbf{Ojas Dubey}}
\institute{Department of Computer Science \& Information Systems\\Birla Institute of Technology and Science, Pilani}
\date{IS F311 Computer Graphics | Final Project Presentation}

\begin{document}

% =========================================================================
% SLIDE 1: Title Slide
% =========================================================================
\begin{frame}
  \titlepage
\end{frame}

% =========================================================================
% SLIDE 2: The Core Problem: Why Direct Volume Rendering?
% =========================================================================
\begin{frame}{The Foundational Problem: Surface Extraction vs. Direct Rendering}
  \begin{columns}[T]
    \begin{column}{0.48\textwidth}
      \begin{block}{The Fallacy of Marching Cubes (1987)}
        \begin{itemize}
          \item \textbf{Binary Surface Ambiguity}: Forces voxels into strict inside/outside states, causing false holes.
          \item \textbf{Partial-Volume Loss}: Blurs sub-voxel biological transitions (e.g., cerebrospinal fluid vs. cortex).
          \item \textbf{Geometric Overhead}: Extracts millions of micro-triangles; overwhelms memory.
        \end{itemize}
      \end{block}
    \end{column}

    \begin{column}{0.48\textwidth}
      \begin{block}{The Direct Volume Rendering Paradigm}
        \begin{itemize}
          \item \textbf{Continuous Optical Medium}: Models scalar fields as light-emitting and absorbing gas.
          \item \textbf{Direct Ray Integration}: Evaluates radiative transport directly from 3D lattice points.
          \item \textbf{Dual 1988 Milestones}:
            \begin{itemize}
              \item \textbf{Marc Levoy}: Continuous isosurfaces.
              \item \textbf{Robert Drebin}: Material mixtures.
            \end{itemize}
        \end{itemize}
      \end{block}
    \end{column}
  \end{columns}

  \vspace{0.4cm}
  \centering
  \small
  \textit{Goal: Render dense bone, semi-transparent soft tissue, and cutaways simultaneously in a single pass.}
\end{frame}

% =========================================================================
% SLIDE 3: Marc Levoy (1988): Continuous Isosurface Display
% =========================================================================
\begin{frame}{Marc Levoy (1988): Continuous Isosurface Formulation}
  \begin{columns}[T]
    \begin{column}{0.52\textwidth}
      \begin{block}{Key Equations \& Principles}
        \begin{itemize}
          \item \textbf{Gradient-Weighted Distance Falloff}:
            $$\alpha(x) = \alpha_v \cdot \max\left(0, \, 1 - \frac{1}{r} \left| \frac{f(x) - f_v}{\|\nabla f(x)\|} \right|\right)$$
            Sharp gradient $\rightarrow$ thin, crisp boundary.\\
            Flat region $\rightarrow$ opacity smoothly drops to zero.
          \item \textbf{Continuous Normal Vectors}:
            $$\vec{N}(x) = \frac{\nabla f(x)}{\|\nabla f(x)\|}, \quad \text{via central differences}$$
          \item \textbf{Multi-Surface Blending}:
            $$\alpha_{\text{total}}(x) = 1 - \prod_{m=1}^{M} (1 - \alpha_m(x))$$
        \end{itemize}
      \end{block}
    \end{column}

    \begin{column}{0.44\textwidth}
      \begin{block}{Core Advantages}
        \begin{itemize}
          \item \textbf{Zero Staircasing}: Continuous trilinear interpolation eliminates voxel stepping.
          \item \textbf{Multi-Tissue Display}: Shows skin envelope ($\alpha_1$) simultaneously with dense cranial bone ($\alpha_2$).
          \item \textbf{Blinn-Phong Shading}: Normal-derived diffuse and specular highlights.
        \end{itemize}
      \end{block}
    \end{column}
  \end{columns}
\end{frame}

% =========================================================================
% SLIDE 4: Robert Drebin (1988): Material Mixture Classification
% =========================================================================
\begin{frame}{Robert Drebin et al. (1988): Probabilistic Material Mixtures}
  \begin{columns}[T]
    \begin{column}{0.52\textwidth}
      \begin{block}{Formulation \& Matte Geometry}
        \begin{itemize}
          \item \textbf{Trapezoidal Material Fractions}:
            $$p_m(f(x)) \in [0, 1], \quad \sum_{m} p_m = 1$$
          \item \textbf{Effective Density \& Surface Strength}:
            $$\rho(x) = \sum_m p_m \rho_m, \quad S(x) = \|\nabla \rho(x)\|$$
          \item \textbf{Interface Accentuation}:
            $$\alpha(x) = \left(\sum_m p_m \alpha_m\right) \cdot \min(1.0, \, k_s S(x))$$
          \item \textbf{3D Matte Cutaways}:
            $$\alpha'(x) = \alpha(x) \cdot M(x)$$
        \end{itemize}
      \end{block}
    \end{column}

    \begin{column}{0.44\textwidth}
      \begin{block}{Visual Capabilities}
        \begin{itemize}
          \item \textbf{Partial Volumes}: Decomposes voxels into air, tissue, and bone percentages.
          \item \textbf{Transparent Interiors}: Homogeneous material interiors ($\|\nabla \rho\| = 0$) stay translucent.
          \item \textbf{Matte Wedge Cuts}: Carves away outer tissues to reveal internal cavities.
        \end{itemize}
      \end{block}
    \end{column}
  \end{columns}
\end{frame}

% =========================================================================
% SLIDE 5: Real Clinical & Scientific Datasets
% =========================================================================
\begin{frame}{Real Scientific \& Clinical Benchmarks}
  \begin{columns}[T]
    \begin{column}{0.48\textwidth}
      \begin{block}{Clinical Neuroimaging}
        \begin{itemize}
          \item \textbf{Human Brain T1 MRI} ($140 \times 160 \times 124$):
            Isotropic $1.0\text{ mm}$ MPRAGE. Sulci/gyri, ventricles, brainstem.
          \item \textbf{Human Brain FLAIR} ($140 \times 160 \times 124$):
            CSF nulled to accentuate tissue boundaries.
          \item \textbf{CT Angiogram AVM} ($337 \times 319 \times 154$):
            Contrast-enhanced cerebral arteriovenous malformation nidus.
        \end{itemize}
      \end{block}
    \end{column}

    \begin{column}{0.48\textwidth}
      \begin{block}{Industrial \& Analytical Benchmarks}
        \begin{itemize}
          \item \textbf{Bonsai Tree Micro-CT} ($256^3 \approx 16.8\text{M voxels}$):
            Stefan R\"ottger. Fine needles, wood rings, roots, clay pot.
          \item \textbf{Engine Block CT} ($256 \times 256 \times 128$):
            GE / Stanford. Specular metal, cylinders, cooling jackets.
          \item \textbf{Marschner-Lobb Benchmark} ($80^3$):
            Analytical cosine-wave field testing reconstruction sharpness.
        \end{itemize}
      \end{block}
    \end{column}
  \end{columns}

  \vspace{0.2cm}
  \centering
  \small
  \textit{All datasets standardized to normalized floating-point ranges $\mathbb{R} \in [0.0, 1.0]$.}
\end{frame}

% =========================================================================
% SLIDE 6: System Architecture: What We Built On Top
% =========================================================================
\begin{frame}{System Architecture: Modern Engineering Additions}
  \begin{columns}[T]
    \begin{column}{0.48\textwidth}
      \begin{block}{Algorithmic Innovations}
        \begin{itemize}
          \item \textbf{Universal Unit-Cube Normalization}:
            Maps arbitrary volumes to $[-1, 1]^3$, permanently resolving the camera-inside-volume bug.
          \item \textbf{Parallel JIT Raycaster (Numba LLVM)}:
            Multi-core work-stealing parallel ray march ($40\times$ faster than interpreted Python).
          \item \textbf{Dual Compositing Pipelines}:
            \begin{itemize}
              \item Back-to-Front (Proposal Specification)
              \item Front-to-Back + Early Ray Termination ($\alpha \ge 0.99 \rightarrow 2.5\times$ speedup)
            \end{itemize}
        \end{itemize}
      \end{block}
    \end{column}

    \begin{column}{0.48\textwidth}
      \begin{block}{Workstation Architecture}
        \begin{itemize}
          \item \textbf{Replaced Tkinter with FastAPI/Canvas}:
            Eliminated single-threaded GUI freezes with async REST server + HTML5 Canvas.
          \item \textbf{Single-Screen 100vh Layout}:
            Zero window scrolling; responsive layout.
          \item \textbf{Progressive Resolution Refinement}:
            Smooth $30\text{ ms}$ orbit previews, crisp $380^2$ frame on release.
          \item \textbf{2D Multi-Planar Orthogonal Slices}:
            Axial, Coronal, and Sagittal cross-section scrubbing.
        \end{itemize}
      \end{block}
    \end{column}
  \end{columns}
\end{frame}

% =========================================================================
% SLIDE 7: Quantitative Results & Conclusion
% =========================================================================
\begin{frame}{Quantitative Benchmarks \& Summary}
  \begin{columns}[T]
    \begin{column}{0.55\textwidth}
      \begin{block}{Performance Across $320 \times 320$ Output}
        \centering
        \footnotesize
        \begin{tabular}{@{}lcccc@{}}
          \toprule
          \textbf{Dataset} & \textbf{Levoy} & \textbf{Levoy (ERT)} & \textbf{Drebin} & \textbf{MIP} \\ \midrule
          Brain MRI & $1.42\text{ s}$ & $\mathbf{0.78\text{ s}}$ & $1.12\text{ s}$ & $0.09\text{ s}$ \\
          Bonsai CT & $1.95\text{ s}$ & $\mathbf{1.08\text{ s}}$ & $1.64\text{ s}$ & $0.12\text{ s}$ \\
          Engine CT & $1.68\text{ s}$ & $\mathbf{0.67\text{ s}}$ & $1.35\text{ s}$ & $0.11\text{ s}$ \\
          AVM Angio & $1.88\text{ s}$ & $\mathbf{0.94\text{ s}}$ & $1.52\text{ s}$ & $0.10\text{ s}$ \\ \bottomrule
        \end{tabular}
      \end{block}

      \begin{block}{Algorithmic Verification}
        \footnotesize
        \textbf{All 6 Unit Tests Passed}: Trilinear interpolation ($10^{-5}$), central difference gradients, slab ray-box intersection, and B2F vs F2B mathematical equivalence ($L_\infty < 10^{-4}$).
      \end{block}
    \end{column}

    \begin{column}{0.42\textwidth}
      \begin{block}{Key Takeaways}
        \begin{itemize}
          \item \textbf{Levoy}: Best for crisp boundary extraction and anatomical surfaces.
          \item \textbf{Drebin}: Best for complex organic mixtures and 3D wedge cutaways.
          \item \textbf{MIP}: Instant preview ($>8\text{ FPS}$) for vascular contrast localization.
          \item \textbf{Workstation}: Live on GitHub at \texttt{TarunPK15/CG-Project}.
        \end{itemize}
      \end{block}
    \end{column}
  \end{columns}
\end{frame}

\end{document}
